\chapter{节奏组：让鼓、贝斯与伴奏咬在同一拍上}
\label{ch:rhythm}

\begin{quote}
\itshape 节奏组不是“鼓加贝斯加一把琴”，而是一台三件套：谁定时间，谁推向前，谁给颜色。
\end{quote}

\section{节奏组到底负责什么}
\label{sec:rhythm-job}

\term{节奏组（rhythm section）}{rhythm-section}是承担律动与根盘的声部群，通常包括鼓、贝斯、节奏吉他/键盘与打击。它至少负责四件事：

\begin{enumerate}
\item \textbf{定时间}：让听众知道一拍有多长、一小节有多长；定得稳，后面的旋律才敢自由。
\item \textbf{定律动}：决定重音落在哪、哪里留空、哪里推进；这是\term{律动（groove）}{groove}的来源。
\item \textbf{定根盘}：用贝斯与底鼓把和声根音与低频重心立住，让上面的和声有地基。
\item \textbf{定能量}：同一首歌里，节奏组的密度与力度决定这一段是“走”还是“冲”。
\end{enumerate}

旋律写得好不好，听众可能记不住；节奏组一乱，听众一定听得出来。所以本章的顺序是：先把时间与根盘写稳，再谈花样。

\section{三层分工：锚点、推进、颜色}
\label{sec:rhythm-layers}

把节奏组拆成三层，比按乐器拆更有用：

\begin{longtable}{p{0.16\textwidth}p{0.30\textwidth}p{0.46\textwidth}}
\toprule
层 & 主要成员 & 职责 \\
\midrule
\endfirsthead
\toprule
层 & 主要成员 & 职责 \\
\midrule
\endhead
锚点 & 底鼓、贝斯根音 & 把每小节的重心与根音钉住；听众靠它判断“拍子在哪”。 \\
推进 & 军鼓、踩镲、贝斯的中段音符 & 把时间从上一拍送到下一拍；决定律动是平、是推、还是切。 \\
颜色 & 打击、节奏吉他/键盘的装饰音 & 给同一律动加不同质感；没有它也能成立，有了它才有风格。 \\
\bottomrule
\end{longtable}

三层不是三组固定乐器。贝斯可以只做锚点，也可以在中段做推进；键盘可以只做颜色，也可以在副歌承担和声。写每一拍之前先问：这一拍是谁在定时间，谁在推，谁在加颜色？三个都想要，就会挤。

\section{鼓：先写骨架，再写密度}
\label{sec:drums}

鼓的写作从三件东西开始：底鼓（kick）、军鼓（snare）、踩镲（hi-hat）。先写它们的落点，再考虑加花与打击。

\begin{longtable}{p{0.20\textwidth}p{0.32\textwidth}p{0.40\textwidth}}
\toprule
骨架 & 落点 & 适用与风险 \\
\midrule
\endfirsthead
\toprule
骨架 & 落点 & 适用与风险 \\
\midrule
\endhead
标准反拍 & 底鼓落 1、3 拍；军鼓落 2、4 拍 & 流行/摇滚最稳的骨架；风险是太普通，需要靠踩镲与加花给性格。 \\
半拍 & 底鼓落 1 拍，军鼓落 3 拍，整体时间感拉宽 & 抒情、慢摇、重氛围；风险是推进不足，副歌容易拖。 \\
四拍四 & 底鼓每拍都落，军鼓在 2、4 或切分 & 舞曲、电子、迪斯科；风险是低频一直顶着，贝斯与底鼓容易打架。 \\
切分底鼓 & 底鼓避开强拍，落在弱拍或反拍 & 放克、嘻哈、拉丁；风险是听众抓不住拍，必须有一个稳定的锚点。 \\
\bottomrule
\end{longtable}

写鼓时遵守两条：

\begin{itemize}
\item \textbf{先有骨架，再有加花。}一小节里底鼓与军鼓的落点必须先能独立成立；加花是装饰，不是骨架。
\item \textbf{军鼓不是律动主角。}军鼓负责反拍与推进，底鼓负责重心；两者抢同一拍，低频与中频会一起糊。
\end{itemize}

踩镲决定密度：主歌可以用四分或八分，副歌再上十六分；不要从头到尾同一密度。打击（percussion）属于颜色层，演示工程默认静音它，就是提醒你：没有它，骨架也要成立。

\section{贝斯：根音优先，律动第二}
\label{sec:bass}

贝斯的第一任务是让听众听到和声的根。第二任务才是让根音“动”起来。顺序反了，贝斯就会变成第二条旋律。

\begin{longtable}{p{0.26\textwidth}p{0.64\textwidth}}
\toprule
写法 & 说明 \\
\midrule
\endfirsthead
\toprule
写法 & 说明 \\
\midrule
\endhead
根音为主 & 每小节或每半小节回到根音；听众靠它确认和弦。 \\
同音重复 & 稳定律动的经验线：同音重复率约 70\% 以上；低于 60\% 容易显得旋律化、抢戏。 \\
控制大跳 & 相邻音的大跳（≥5 半音）占比控制在 10\% 以内；经过句可以破例，但不要整段都在跳。 \\
音区 & 常见起点 E1—G2（28—43）；频繁上探到 C3 以上，低频重心会空，听感也更容易“飘”。 \\
时值 & 长音更稳，八分更有推进；主歌用长一点，副歌再切短，不要一开始就把贝斯写成十六分。 \\
\bottomrule
\end{longtable}

贝斯与底鼓的关系，比贝斯与军鼓更重要。两种基本对齐：

\begin{itemize}
\item \textbf{同落}：贝斯与底鼓同时落重拍，低频最稳；适合大多数流行/摇滚的根盘。
\item \textbf{错落}：底鼓落强拍，贝斯晚半拍或落反拍，律动更有推动；适合放克、电子，但需要更强的锚点防止散。
\end{itemize}

\section{伴奏：压缩和弦，不要每拍都满}
\label{sec:comping}

节奏吉他/键盘的\term{压缩伴奏（comping）}{comping}不是“把和弦从头按到尾”。它的工作是：在正确的拍点给和弦，在不需要的地方留空。

\begin{itemize}
\item \textbf{给和弦的骨架}：先确定每小节哪几拍必须听到和弦变化；其余拍可以留空或只给单音。
\item \textbf{避开贝斯根音}：伴奏的低音声部不要和贝斯抢同一个根音；把和弦放在中频，低频留给贝斯。
\item \textbf{每四小节留一点空}：持续的八分或十六分会把听众听累；留一拍真空，鼓与贝斯的律动会更清楚。
\item \textbf{扫弦与单音二选一}：满扫弦适合大开大合；单音 riff 更适合沉、稳的律动。不要每一拍都又扫又填。
\end{itemize}

\section{三轨对齐：一张网格上的三个问题}
\label{sec:rhythm-align}

把鼓、贝斯、伴奏放在同一张八拍网格上，只需要检查三个问题：

\begin{enumerate}
\item \textbf{底鼓与贝斯是否咬住？}同落时低频是否太重，错落时是否有锚点？
\item \textbf{军鼓与伴奏是否互相遮？}军鼓落 2、4 拍时，伴奏不要在同一拍给满和弦。
\item \textbf{踩镲与打击是否太密？}密度是相对段落而言的；主歌的十六分，到了副歌可能只是基础。
\end{enumerate}

\begin{figure}[htbp]
  \centering
  \begin{tikzpicture}[font=\small]
    \draw[step=1.2cm, gray!35, thin] (0,0) grid (9.6,3.6);
    \draw[thick] (0,0) rectangle (9.6,3.6);
    \node[anchor=east] at (-0.1,3.0) {底鼓};
    \node[anchor=east] at (-0.1,2.1) {军鼓};
    \node[anchor=east] at (-0.1,1.2) {贝斯};
    \node[anchor=east] at (-0.1,0.3) {伴奏};
    \foreach \x in {0.1,2.5,4.9,7.3} {
      \fill[yebanblue!75] (\x,2.75) rectangle ++(1.2,0.5);
    }
    \foreach \x in {1.3,3.7,6.1,8.5} {
      \fill[gray!70] (\x,1.85) rectangle ++(1.2,0.5);
    }
    \foreach \x in {0.1,1.3,2.5,3.7,4.9,6.1,7.3,8.5} {
      \fill[yebanblue!30] (\x,0.95) rectangle ++(1.1,0.45);
    }
    \foreach \x in {0.1,2.5,4.9,7.3} {
      \fill[gray!35] (\x,0.05) rectangle ++(1.2,0.45);
    }
    \node[font=\scriptsize, anchor=west] at (9.7,3.0) {强拍};
    \node[font=\scriptsize, anchor=west] at (9.7,2.1) {反拍};
    \node[font=\scriptsize, anchor=west] at (9.7,1.2) {根音};
    \node[font=\scriptsize, anchor=west] at (9.7,0.3) {留空};
  \end{tikzpicture}
  \caption{八拍网格上的节奏组示意：底鼓定强拍，军鼓落反拍，贝斯守根音，伴奏在关键拍给和弦并在其余拍留空。}
  \label{fig:rhythm-grid}
\end{figure}

这张图不是让你照抄落点，而是让你在写每一拍时能指出“这一拍谁负责什么”。如果四个行都在同一拍落满，听众听到的不是丰富，而是挤。

\section{在夜半里落地}
\label{sec:rhythm-yeban}

\begin{enumerate}
\item \textbf{先看清演示工程的角色。}演示工程的鼓是 MIDI 轨，贝斯是音频轨，打击是默认静音的音频轨。听成品时可以 Solo 贝斯与鼓，观察它们怎样咬合；自己写贝斯时建议新建 MIDI 合成器轨，才能改音符、力度与时值。
\item \textbf{建轨。}「轨道 → + 新建 MIDI 合成器轨（Ctrl+N）」建鼓/贝斯/伴奏的 MIDI 轨；音频素材放音频轨。不要把鼓、贝斯、伴奏全塞进同一条轨。
\item \textbf{在钢琴卷帘里写骨架。}先写底鼓与军鼓的落点，再写贝斯的根音，最后写伴奏和弦；顺序反过来，你会被装饰音带着走。
\item \textbf{量化与人性化。}用「编辑 → 音符量化（Quantize 1/16）」把手弹的毛边对齐网格；再用「编辑 → 音符人性化（Humanize）」给时值与力度可复现的微差，不要把全曲量化成机械直线。
\item \textbf{力度泳道。}在钢琴卷帘下方逐个画力度；律动声部不要全曲同一个力度。主歌 90—105、副歌 105—120 是常见起点，不是硬规定。
\item \textbf{Solo/Mute 检查。}Solo 鼓＋贝斯听低频是否咬住；Mute 伴奏听骨架是否还成立；Mute 打击听颜色层是否只是装饰。
\item \textbf{混音台看低频。}按 X / F8 打开混音台，观察底鼓与贝斯的电平是否互相顶；编曲阶段先解决“同一拍抢低频”，不要用 EQ 硬切。
\end{enumerate}

\screenshot{arrangement-rhythm-section-groove}{鼓、贝斯、伴奏三条轨的八拍骨架与混音台低频电平：看清谁定强拍、谁落反拍、谁在留空}

\section{常见错误与修正}
\label{sec:rhythm-mistakes}

\begin{longtable}{p{0.34\textwidth}p{0.58\textwidth}}
\toprule
常见错误 & 修正办法 \\
\midrule
\endfirsthead
\toprule
常见错误 & 修正办法 \\
\midrule
\endhead
贝斯写成第二条旋律 & 回到根音；同音重复率提上去，大跳压到 10\% 以内，先让听众听到和声。 \\
底鼓与贝斯每拍都同落 & 改成一强一弱或错半拍；低频一直双顶，听感会轰。 \\
军鼓与伴奏同一拍给满 & 伴奏在军鼓落点留空或只给单音；把反拍让给军鼓。 \\
踩镲从头到尾十六分 & 按段落分密度：主歌八分、副歌十六分；用留空制造推进。 \\
全曲量化得一样整齐 & 骨架对齐网格，装饰音与力度做人性化；不要把所有轨都量化成同一格。 \\
力度全曲一条直线 & 主歌/副歌各留动态层级；律动声部也要有呼吸。 \\
伴奏低音与贝斯抢根音 & 伴奏把低音声部撤掉或升高八度，低频留给贝斯。 \\
打击从头铺到尾 & 打击是颜色层；没有它骨架也要成立，按段落加减。 \\
\bottomrule
\end{longtable}

\section{练习}
\label{sec:rhythm-exercise}

\begin{exercise}
做一条八小节节奏组，完成三版并记录差异：
\begin{enumerate}
\item \textbf{标准版}：底鼓 1、3，军鼓 2、4，贝斯根音为主，伴奏每两拍给一次和弦。
\item \textbf{推进版}：保持骨架，副歌把踩镲加到十六分，贝斯在每小节第三拍加一个经过音；记录同音重复率与大跳占比。
\item \textbf{半拍版}：底鼓落 1，军鼓落 3，贝斯时值拉长；比较三版的低频重心与“走/拖”的感觉。
\end{enumerate}
每版都用 Solo/Mute 检查：Solo 鼓＋贝斯时骨架是否成立；Mute 伴奏时律动是否还在；Mute 打击时颜色是否只是装饰。
\end{exercise}

\section{本章要点}
\label{sec:rhythm-takeaways}

\begin{itemize}
\item 节奏组负责定时间、定律动、定根盘、定能量。
\item 三层分工：锚点、推进、颜色；每一拍先问谁负责什么。
\item 鼓先写骨架再写密度；军鼓不是律动主角。
\item 贝斯根音优先；同音重复率约 70\% 以上、大跳占比 10\% 以内是稳定律动的经验线。
\item 伴奏要压缩和弦、留空、避开贝斯根音；不要每拍都满。
\item 编曲阶段先解决低频同拍冲突，不要留给 EQ 与压缩救火。
\end{itemize}

\begin{helpmapping}
\begin{longtable}{p{0.45\textwidth}p{0.45\textwidth}}
\toprule
本节 & 帮助锚点 \\
\midrule
节奏组到底负责什么 & \texttt{arrangement-rhythm-job} \\
三层分工 & \texttt{arrangement-rhythm-layers} \\
鼓：先写骨架，再写密度 & \texttt{arrangement-drums} \\
贝斯：根音优先 & \texttt{arrangement-bass} \\
伴奏：压缩和弦 & \texttt{arrangement-comping} \\
三轨对齐 & \texttt{arrangement-rhythm-align} \\
在夜半里落地 & \texttt{arrangement-rhythm-yeban} \\
常见错误与修正 & \texttt{arrangement-rhythm-mistakes} \\
练习 & \texttt{arrangement-rhythm-exercise} \\
\bottomrule
\end{longtable}
\end{helpmapping}
